\section{Background and Scope}

\subsection{Video Prediction as Predictive Modeling}

Video prediction studies how a model can infer future visual observations from past visual evidence, sometimes with additional conditioning variables such as actions. In its simplest form, the model receives one or more frames and predicts subsequent frames; in action-conditioned settings, the prediction also depends on the intervention an agent will take. Early action-conditioned work made this connection explicit by predicting high-dimensional Atari observations under discrete control inputs, while robot-interaction video prediction used action-conditioned motion transformations to forecast how objects move under manipulation \cite{oh2015_action_conditional_video_prediction_prov,finn2016_physical_interaction_video_prediction_prov}. SV2P illustrates why this task cannot always be reduced to a single deterministic future: latent variables can represent multiple plausible rollouts when the future is ambiguous \cite{babaeizadeh2018_sv2p_prov}.

\subsection{World Models Beyond Future Frames}

In this survey, a world model is a learned predictive model that represents how an environment evolves and can support controlled prediction, planning, policy learning, or simulator-style analysis. This usage is grounded in latent world-model work, where images are compressed into compact states and dynamics are learned over those states. The recurrent world-model formulation factorizes perception, memory, and control into VAE, MDN-RNN, and controller modules; PlaNet plans directly in a recurrent latent state; and Dreamer trains behavior through imagined latent trajectories rather than relying only on pixel-space prediction \cite{ha2018_world_models_prov,hafner2019_planet_prov,hafner2020_dreamer_prov}.

These examples illustrate why predictive world models from videos are broader than future-frame prediction: the learned model may be evaluated by control return, planning success, or the quality of imagined interaction, not only by frame reconstruction. A frame-level model can be useful for control even when pixel metrics do not fully capture its value, and an object-structured representation can support planning better than an object-agnostic video predictor in a structured physical task \cite{oh2015_action_conditional_video_prediction_prov,janner2019_o2p2_prov}. DIAMOND adds a related evaluation case in which an agent is trained inside an image-space diffusion world model rather than only judging generated frames \cite{alonso2024_diamond_prov}.
Recent feature-space models such as V-JEPA 2 further motivate this boundary because they are discussed as predictive video representations with reported planning use, shifting attention from frame reconstruction toward behavior-oriented evaluation \cite{assran2025vjepa2_prov}.

This broader view also changes how the survey treats visual quality. VideoPhy suggests that visual plausibility should not be equated with accurate physical simulation, because generated videos can remain inconsistent with the physical interactions implied by prompts \cite{bansal2025_videophy_prov}. The survey therefore treats visual realism, predictive accuracy, controllability, and physical consistency as related but distinct criteria.

\subsection{Axes for Organizing the Survey}

The first organizing axis is representation. Pixel-level models directly predict or transform image observations, as in action-conditioned Atari prediction and robot pushing \cite{oh2015_action_conditional_video_prediction_prov,finn2016_physical_interaction_video_prediction_prov}. Latent world models instead learn compact predictive states for planning or policy learning \cite{ha2018_world_models_prov,hafner2019_planet_prov,hafner2020_dreamer_prov}. Object-centric models represent entities and relations, using graph-based or object-factorized dynamics to model physical interactions \cite{battaglia2016_interaction_networks_prov,watters2017_visual_interaction_networks_prov,janner2019_o2p2_prov}. DIAMOND exemplifies an image-space diffusion world model that predicts observations from action and observation history \cite{alonso2024_diamond_prov}. Interactive generative environments add another representational layer by learning latent action spaces from video-only data \cite{bruce2024_genie_prov}.

The second organizing axis is dynamics. Pixel-space predictors include autoregressive, recurrent, motion-transform, and stochastic latent mechanisms for rolling visual observations forward \cite{finn2016_physical_interaction_video_prediction_prov,babaeizadeh2018_sv2p_prov}. Latent control models learn recurrent state-space dynamics over compact states, with PlaNet using latent-space planning and Dreamer using actor-critic learning through imagined rollouts \cite{hafner2019_planet_prov,hafner2020_dreamer_prov}. Object-centric models instead place dynamics on objects and relations, either assuming an explicit graph or learning a visual front end that feeds a relational predictor \cite{battaglia2016_interaction_networks_prov,watters2017_visual_interaction_networks_prov}. These alternatives make the dynamics model a central taxonomy axis rather than an implementation detail.

The third organizing axis is conditioning and interaction. A passive video predictor extrapolates from observed frames, whereas a world model used for control requires some account of interventions. Actions enter explicitly in Atari video prediction, robot pushing, and latent control models \cite{oh2015_action_conditional_video_prediction_prov,finn2016_physical_interaction_video_prediction_prov,hafner2019_planet_prov,hafner2020_dreamer_prov}. Genie broadens this picture by inferring discrete latent actions from unlabeled video instead of relying on ground-truth control labels \cite{bruce2024_genie_prov}. This distinction matters because world models are often used not merely to forecast what will happen, but to evaluate what could happen under alternative actions.

The fourth organizing axis is evaluation. Frame metrics such as reconstruction error, PSNR, and SSIM are useful for some video-prediction settings, but they do not by themselves establish planning utility or physical correctness. Action-conditioned prediction can be evaluated through control usefulness, object-oriented prediction can be evaluated through downstream planning, and diffusion-based world models can be evaluated through agent performance inside the learned model \cite{oh2015_action_conditional_video_prediction_prov,janner2019_o2p2_prov,alonso2024_diamond_prov}. Physical-consistency benchmarks add another layer by asking whether generated videos obey physical commonsense, with VideoPhy organizing prompts around solid-solid, solid-fluid, and fluid-fluid interactions and using both human and automatic evaluation \cite{bansal2025_videophy_prov}. This survey therefore tracks which aspect of world modeling is being tested: visual fidelity, predictive rollout, control utility, interaction, physical consistency, or human-aligned judgment.

\subsection{Survey Scope and Conceptual Families}

This survey distinguishes six conceptual families. Pixel-level video prediction models future visual observations directly and provides one historical bridge from passive sequence modeling to action-conditioned prediction \cite{oh2015_action_conditional_video_prediction_prov,finn2016_physical_interaction_video_prediction_prov,babaeizadeh2018_sv2p_prov}. Latent world models compress video observations into predictive states for planning or learned behavior \cite{ha2018_world_models_prov,hafner2019_planet_prov,hafner2020_dreamer_prov}. Object-centric physical dynamics models emphasize entities, relations, hidden physical variables, and planning-relevant structure \cite{battaglia2016_interaction_networks_prov,watters2017_visual_interaction_networks_prov,janner2019_o2p2_prov}. DIAMOND anchors the diffusion-based world-model family considered here, where image-space generative dynamics raise the question of when visual detail should be preserved rather than abstracted away \cite{alonso2024_diamond_prov}. Interactive and embodied world models condition predictions on actions or learned latent controls, shifting the goal from watching video to generating controllable futures \cite{oh2015_action_conditional_video_prediction_prov,finn2016_physical_interaction_video_prediction_prov,bruce2024_genie_prov}. Finally, evaluation benchmarks for physical consistency test whether visually plausible videos also satisfy physical commonsense and interaction constraints \cite{bansal2025_videophy_prov}.

The next section formalizes these distinctions into a taxonomy of modeling goals, representations, dynamics, conditioning signals, and evaluation protocols.
